\section{Descomposición en valores singulares (SVD)}

\begin{teoria}
La SVD generaliza la diagonalización ortogonal a matrices \emph{no cuadradas}. Sea $A\in\R^{n\times m}$. Existen $U\in\R^{n\times n}$ ortogonal ($U^TU=UU^T=I$), $V\in\R^{m\times m}$ ortogonal y $\Sigma\in\R^{n\times m}$ ``diagonal rectangular'' tales que
\[
A = U\,\Sigma\,V^T,
\qquad
\Sigma=\diag(\sigma_1,\dots,\sigma_r)\ \text{orlada con ceros hasta tamaño } n\times m,
\qquad
\sigma_1\ge\sigma_2\ge\dots\ge\sigma_r>0.
\]
A los $\sigma_i$ se los llama \textbf{valores singulares} de $A$. La forma exacta de $\Sigma$ depende de $n$ frente a $m$: lleva los $\sigma_i$ en la diagonal y se completa con ceros. Si se trabaja en $\C$, la descomposición es $A=U\Sigma V^*$ con $V^*=\overline{V}^{\,T}$ (conjugada transpuesta) y $U$ unitaria.

\medskip
\textbf{De dónde salen las piezas.} Como $U,V$ son ortogonales,
\[
A^TA = (U\Sigma V^T)^T(U\Sigma V^T) = V\,\Sigma^T\Sigma\,V^T,
\qquad
AA^T = U\,\Sigma\Sigma^T\,U^T .
\]
Es decir: $V$ diagonaliza ortogonalmente a $A^TA$ y $U$ diagonaliza ortogonalmente a $AA^T$. Los autovalores no nulos de $A^TA$ y de $AA^T$ coinciden (cambian sólo las multiplicidades de $0$), y
\[
\sigma_i=\sqrt{\lambda_i(A^TA)}\,.
\]

\textbf{Relaciones fundamentales} (con $r=\rango(A)\le\min(n,m)$): existen bases ortonormales $\{u_1,\dots,u_n\}$ de $\R^n$ y $\{v_1,\dots,v_m\}$ de $\R^m$ tales que
\begin{align*}
A v_i &= \sigma_i u_i, & A^T u_i &= \sigma_i v_i, & i&=1,\dots,r,\\
A v_i &= 0, & A^T u_i &= 0, & i&>r.
\end{align*}
De aquí sale el atajo $u_i=\dfrac{Av_i}{\sigma_i}$ (o $v_i=\dfrac{A^Tu_i}{\sigma_i}$), que evita calcular el segundo producto matricial.

\textbf{Rango y número de condición.} El número de valores singulares no nulos es exactamente el rango:
\[
\rango(A)=\#\{\sigma_i>0\},
\qquad
\#\{\sigma_i=0\}=\min(n,m)-\rango(A).
\]
El \textbf{número de condición} en norma $2$ es
\[
\kappa_2(A)=\frac{\sigma_{\max}}{\sigma_{\min}} .
\]
Cuando $\sigma_{\min}\to 0$ (matriz casi singular) $\kappa_2\to\infty$: el mal condicionamiento se lee en el valor singular más chico. Si $A$ es simétrica/hermítica, $\sigma_i=\abs{\lambda_i}$ y $\kappa_2(A)=\dfrac{\max_i\abs{\lambda_i}}{\min_i\abs{\lambda_i}}$.
\end{teoria}

\begin{receta}
\textbf{Cálculo de la SVD $A=U\Sigma V^T$ (vía $A^TA$).}
\begin{enumerate}
\item Formar $A^TA$ (matriz $m\times m$, simétrica). Hallar sus autovalores $\lambda_1\ge\dots\ge\lambda_m\ge0$.
\item Valores singulares: $\sigma_i=\sqrt{\lambda_i}$. Armar $\Sigma\in\R^{n\times m}$ con los $\sigma_i$ ordenados de mayor a menor en la diagonal.
\item Autovectores de $A^TA$ normalizados $\to$ columnas de $V$. Ortonormalizar dentro de cada autoespacio si hay repeticiones.
\item Para $\sigma_i>0$: $u_i=\dfrac{Av_i}{\sigma_i}$. Completar $u_{r+1},\dots,u_n$ a una base ortonormal de $\R^n$ (por ejemplo con Gram--Schmidt o producto vectorial).
\end{enumerate}
\textbf{Vía $AA^T$ (simétrica):} análogo intercambiando roles: se obtiene $U$ de $AA^T$ y luego $v_i=\dfrac{A^Tu_i}{\sigma_i}$, completando $V$ con una base de $\Nuc(A)$.
\end{receta}

\begin{ejemplo}{Guía VII, Ej.\ 1 (SVD de una $4\times3$)}
Hallar la SVD de
\[
A=\begin{pmatrix} 0 & 3 & 3 \\ 4 & 1 & -1 \\ 4 & 1 & -1 \\ 0 & 3 & 3 \end{pmatrix}.
\]
Como $AA^T$ es $4\times4$ pero da bloques simples, se puede trabajar con $A^TA$ ($3\times3$):
\[
A^TA=\begin{pmatrix} 32 & 8 & -8 \\ 8 & 20 & 16 \\ -8 & 16 & 20 \end{pmatrix},
\qquad
p(\lambda)=-\lambda(\lambda-36)^2 .
\]
Entonces $\lambda_1=\lambda_2=36$, $\lambda_3=0$, de donde $\sigma_1=\sigma_2=6$, $\sigma_3=0$ y
\[
\Sigma=\begin{pmatrix} 6 & 0 & 0 \\ 0 & 6 & 0 \\ 0 & 0 & 0 \\ 0 & 0 & 0 \end{pmatrix}.
\]
Autovectores de $A^TA$ (autoespacio de $\lambda=36$, ya ortogonalizado): $v_1=\tfrac13(2,2,1)^T$, $v_2=\tfrac13(2,-1,-2)^T$. Luego $u_i=\tfrac1{6}Av_i$ da
\[
u_1=\tfrac12(1,1,1,1)^T,\qquad u_2=\tfrac12(-1,1,1,-1)^T,
\]
y se completa la base de $\R^4$ con $(1,0,0,-1)^T$, $(0,1,-1,0)^T$ (ortogonales a $u_1,u_2$). Forma reducida:
\[
A=\tfrac12\begin{pmatrix} 1 & -1 \\ 1 & 1 \\ 1 & 1 \\ 1 & -1 \end{pmatrix}
\begin{pmatrix} 6 & 0 \\ 0 & 6 \end{pmatrix}
\tfrac13\begin{pmatrix} 2 & 2 & 1 \\ 2 & -1 & -2 \end{pmatrix}.
\]
Como $\sigma_3=0$, $\rango(A)=2$.
\end{ejemplo}

\begin{ejemplo}{Final Tema VIII, Ej.\ 3 (SVD de una $2\times3$)}
Hallar una SVD de $A=\begin{pmatrix} 1 & -3 & 2 \\ 5 & 2 & 0 \end{pmatrix}$. Al ser $2\times3$ conviene $AA^T$ ($2\times2$):
\[
AA^T=\begin{pmatrix} 14 & -1 \\ -1 & 29 \end{pmatrix},
\qquad
p(\lambda)=\lambda^2-43\lambda+405,
\qquad
\lambda_{1,2}=\frac{43\pm\sqrt{229}}{2}.
\]
Valores singulares (orden decreciente):
\[
\sigma_1=\sqrt{\tfrac{43+\sqrt{229}}{2}}\approx 5.3913,
\qquad
\sigma_2=\sqrt{\tfrac{43-\sqrt{229}}{2}}\approx 3.7328 .
\]
Chequeo: $\lambda_1\lambda_2=405=\det(AA^T)$ y $\lambda_1+\lambda_2=43=\tr(AA^T)$. Los autovectores de $AA^T$ dan las columnas de $U$; luego $v_i=\dfrac{A^Tu_i}{\sigma_i}$ para $i=1,2$ y $v_3$ es una base ortonormal de $\Nuc(A)$ (resolviendo $Ax=0$ se obtiene $\tilde v_3=(-4,10,17)^T$, $\norm{\tilde v_3}=9\sqrt5$). Resultado numérico:
\[
U\approx\begin{pmatrix} 0.0662 & 0.9978 \\ -0.9978 & 0.0662 \end{pmatrix},
\quad
\Sigma\approx\begin{pmatrix} 5.3913 & 0 & 0 \\ 0 & 3.7328 & 0 \end{pmatrix},
\quad
V\approx\begin{pmatrix} -0.9131 & 0.3560 & -0.1988 \\ -0.4070 & -0.7664 & 0.4969 \\ 0.0246 & 0.5346 & 0.8447 \end{pmatrix}.
\]
Se verifica $U\Sigma V^T=A$, $U^TU=I_2$, $V^TV=I_3$.
\end{ejemplo}

\begin{truco}
\begin{itemize}
\item \textbf{Ordená siempre $\sigma_1\ge\sigma_2\ge\dots\ge0$}. Si desordenás los $\sigma_i$, las columnas de $U$ y $V$ dejan de emparejarse y no reconstruís $A$.
\item \textbf{$A^TA$ vs.\ $AA^T$}: elegí el producto \emph{más chico}. Para $A\in\R^{n\times m}$ con $n<m$ (ancha) calculá $AA^T$ ($n\times n$); si $n>m$ (alta) calculá $A^TA$ ($m\times m$). Los autovalores no nulos son los mismos; sólo cambian las multiplicidades del $0$ y qué matriz ($U$ o $V$) sale directa.
\item \textbf{$\#\{\sigma_i=0\}=\min(n,m)-\rango(A)$}: para saber cuántos valores singulares nulos hay \emph{no hace falta calcular la SVD}, basta el rango. Si una fila/columna es combinación lineal de otras, hay al menos un $\sigma_i=0$.
\item Verificá con las trazas y determinantes: $\sum\lambda_i=\tr(A^TA)$, $\prod\lambda_i=\det(A^TA)$.
\end{itemize}
\end{truco}


\section{Pseudoinversa de Moore--Penrose}

\begin{teoria}
La \textbf{pseudoinversa de Moore--Penrose} $A^{+}$ generaliza la inversa a matrices rectangulares o singulares. A partir de la SVD $A=U\Sigma V^T$ se define
\[
A^{+}=V\,\Sigma^{+}\,U^T,
\]
donde $\Sigma^{+}\in\R^{m\times n}$ se obtiene \emph{transponiendo} $\Sigma$ e invirtiendo los valores singulares no nulos:
\[
\text{si }\Sigma=\diag(\sigma_1,\dots,\sigma_r)\ \text{(orlada con ceros)}
\quad\Longrightarrow\quad
\Sigma^{+}=\diag(1/\sigma_1,\dots,1/\sigma_r)^{T}\ \text{(orlada con ceros)}.
\]
Es decir: cada $\sigma_i>0$ pasa a $1/\sigma_i$, los ceros quedan en cero, y se transpone la forma rectangular.

\medskip
\textbf{Relación con las ecuaciones normales.} Si $\det(A^TA)\ne0$ (columnas de $A$ linealmente independientes),
\[
A^{+}=(A^TA)^{-1}A^T
\qquad\text{(caso sobredeterminado, $n>m$)},
\]
y esta $A^+$ es la que aparece en cuadrados mínimos: $x=A^{+}b$ resuelve $A^TAx=A^Tb$. Análogamente, si las filas son LI ($n<m$, caso subdeterminado) $A^{+}=A^T(AA^T)^{-1}$ y $x=A^+b$ es la solución de norma mínima de $Ax=b$. Si $A$ es cuadrada e inversible, $A^{+}=A^{-1}$.

\textbf{Propiedades (Penrose).} $A^+$ es la única matriz que cumple
\[
AA^{+}A=A,\qquad A^{+}AA^{+}=A^{+},\qquad (AA^{+})^T=AA^{+},\qquad (A^{+}A)^T=A^{+}A .
\]
En particular, si $\rango(A)=m$ (columnas LI) entonces $A^{+}A=I_m$; si $\rango(A)=n$ (filas LI) entonces $AA^{+}=I_n$.
\end{teoria}

\begin{receta}
\textbf{Pseudoinversa por SVD.}
\begin{enumerate}
\item Calcular la SVD $A=U\Sigma V^T$ (receta de la sección anterior).
\item Formar $\Sigma^{+}$: transponer $\Sigma$ y reemplazar cada $\sigma_i>0$ por $1/\sigma_i$ (los ceros quedan).
\item $A^{+}=V\Sigma^{+}U^T$.
\end{enumerate}
\textbf{Atajo sin SVD} (si $A$ tiene columnas LI): $A^{+}=(A^TA)^{-1}A^T$.
\end{receta}

\begin{ejemplo}{Pizarrón SVD / Guía VII Ej.\ 2 (pseudoinversa de una $2\times3$)}
Sea $C=\begin{pmatrix} 3 & 2 & 2 \\ 2 & 3 & -2 \end{pmatrix}$. De la SVD (calculada vía $CC^T=\begin{pmatrix}17&8\\8&17\end{pmatrix}$, con $\lambda=25,9$, $\sigma_1=5$, $\sigma_2=3$) se obtienen
\[
U=\tfrac1{\sqrt2}\begin{pmatrix} 1 & 1 \\ 1 & -1 \end{pmatrix},
\quad
\Sigma=\begin{pmatrix} 5 & 0 & 0 \\ 0 & 3 & 0 \end{pmatrix},
\quad
V=\begin{pmatrix} 1/\sqrt2 & 1/(3\sqrt2) & 2/3 \\ 1/\sqrt2 & -1/(3\sqrt2) & -2/3 \\ 0 & 4/(3\sqrt2) & -1/3 \end{pmatrix}.
\]
La pseudoinversa de $\Sigma$ es
\[
\Sigma^{+}=\begin{pmatrix} 1/5 & 0 \\ 0 & 1/3 \\ 0 & 0 \end{pmatrix},
\]
y entonces
\[
C^{+}=V\,\Sigma^{+}\,U^T=\frac{1}{45}\begin{pmatrix} 7 & 2 \\ 2 & 7 \\ 10 & -10 \end{pmatrix}.
\]
\end{ejemplo}

\begin{ejemplo}{Pizarrón SVD (pseudoinversa de un vector columna)}
Sea $A=\begin{pmatrix}1\\3\end{pmatrix}\in\R^{2\times1}$. Entonces $A^TA=10$, $\sigma_1=\sqrt{10}$, $V=(1)$ y
\[
U=\tfrac1{\sqrt{10}}\begin{pmatrix} 1 & -3 \\ 3 & 1 \end{pmatrix},
\qquad
\Sigma=\begin{pmatrix}\sqrt{10}\\0\end{pmatrix}.
\]
Con $\Sigma^{+}=\big(\tfrac1{\sqrt{10}}\ \ 0\big)$ resulta
\[
A^{+}=V\Sigma^{+}U^T=\tfrac{1}{10}\begin{pmatrix}1 & 3\end{pmatrix},
\qquad A^{+}A=1 .
\]
(Coincide con $A^+=(A^TA)^{-1}A^T=\tfrac1{10}(1\ 3)$.)
\end{ejemplo}

\begin{truco}
\begin{itemize}
\item $\Sigma^{+}$ \textbf{se transpone}: si $\Sigma$ es $n\times m$, entonces $\Sigma^{+}$ es $m\times n$. Un error típico es olvidar transponer la forma rectangular.
\item Nunca inviertas un $\sigma_i=0$: los ceros de $\Sigma$ quedan como ceros en $\Sigma^{+}$.
\item Chequeo rápido de dimensiones: $A^{+}$ es $m\times n$ (transpuesta de la forma de $A$). Verificá con $AA^{+}A=A$.
\item Si sólo necesitás resolver cuadrados mínimos y las columnas son LI, usá directamente $(A^TA)^{-1}A^T$ sin pasar por la SVD completa.
\end{itemize}
\end{truco}


\section{Cuadrados mínimos}

\begin{teoria}
Dado un sistema \emph{sobredeterminado} $Ax=b$ ($A\in\R^{n\times m}$, $n>m$, sin solución exacta), el problema de \textbf{cuadrados mínimos} busca
\[
\min_{x}\ \norm{Ax-b}_2^2 = \min_x\ \inner{Ax-b}{Ax-b}.
\]
El minimizador satisface las \textbf{ecuaciones normales}
\[
A^TA\,x = A^Tb .
\]
Geométricamente: $Ax^\star$ es la proyección ortogonal de $b$ sobre $\Imag(A)$, y el residuo $b-Ax^\star$ es ortogonal a las columnas de $A$.

\medskip
\textbf{Tres vías de resolución.}
\begin{itemize}
\item \textbf{Ecuaciones normales / Cholesky.} Si las columnas de $A$ son LI, $A^TA$ es simétrica definida positiva, luego admite Cholesky $A^TA=LL^T$. Se resuelve $Lz=A^Tb$ (sustitución hacia adelante) y $L^Tx=z$ (hacia atrás). Barato pero empeora el condicionamiento ($\kappa_2(A^TA)=\kappa_2(A)^2$).
\item \textbf{QR.} Con $A=QR$ ($Q$ de columnas ortonormales, $R$ triangular superior), las ecuaciones normales se reducen a $Rx=Q^Tb$. Numéricamente más estable que Cholesky.
\item \textbf{SVD / pseudoinversa.} La solución es $x=A^{+}b=V\Sigma^{+}U^Tb$. Es la vía más robusta y funciona incluso con columnas linealmente dependientes (entrega la solución de norma mínima).
\end{itemize}

\textbf{Ajuste de curvas.} La clave de examen: un modelo puede ser \emph{no lineal en la variable} pero \textbf{lineal en los parámetros}, y entonces es un problema de cuadrados mínimos lineal. Para un modelo $y\approx\sum_j a_j\,\phi_j(x)$ con funciones base $\phi_j$ (potencias, senos, cosenos, etc.), se arma la matriz de diseño con una columna por cada $\phi_j$ evaluada en los $x_k$, y se resuelve por ecuaciones normales.
\end{teoria}

\begin{receta}
\textbf{Ajuste por cuadrados mínimos (modelo lineal en los parámetros).}
\begin{enumerate}
\item Elegir el modelo $y\approx a_1\phi_1(x)+\dots+a_m\phi_m(x)$.
\item Armar la matriz de diseño $A$ (fila $k$ = $(\phi_1(x_k),\dots,\phi_m(x_k))$) y el vector $b=(y_1,\dots,y_n)^T$.
\item Formar $A^TA$ y $A^Tb$.
\item Resolver $A^TAx=A^Tb$ por Cholesky (o QR, o $x=A^{+}b$ por SVD).
\end{enumerate}
\end{receta}

\begin{ejemplo}{Final 2019 (ajuste polinómico $y=ax^3+bx+c$ por Cholesky)}
Ajustar $y=ax^3+bx+c$ a la tabla $(x_k,y_k)$: $(-2,-9),(-1,0),(0,3),(1,6),(2,15)$. Columnas $[x^3,\ x,\ 1]$:
\[
A=\begin{pmatrix} -8 & -2 & 1 \\ -1 & -1 & 1 \\ 0 & 0 & 1 \\ 1 & 1 & 1 \\ 8 & 2 & 1 \end{pmatrix},
\qquad
b=\begin{pmatrix} -9 \\ 0 \\ 3 \\ 6 \\ 15 \end{pmatrix}.
\]
Ecuaciones normales $A^TA\,v=A^Tb$ con $v=(a,b,c)^T$:
\[
A^TA=\begin{pmatrix} 130 & 34 & 0 \\ 34 & 10 & 0 \\ 0 & 0 & 5 \end{pmatrix},
\qquad
A^Tb=\begin{pmatrix} 198 \\ 54 \\ 15 \end{pmatrix}.
\]
Cholesky $A^TA=LL^T$ (el bloque $2\times2$): $\ell_{11}=\sqrt{130}$, $\ell_{21}=34/\sqrt{130}$, $\ell_{22}=\sqrt{10-34^2/130}\approx1.0525$; el bloque $(3,3)$ da $\sqrt5$. Resolviendo $Lz=A^Tb$ y $L^Tv=z$:
\[
a=1,\qquad b=2,\qquad c=3
\qquad\Longrightarrow\qquad y=x^3+2x+3 .
\]
Los datos son exactos: el residuo es $0$ (por ejemplo $x=-2\to-9$, $x=2\to15$).
\end{ejemplo}

\begin{ejemplo}{Final 2009 (ajuste de temperatura, modelo lineal en los parámetros)}
Se modela la temperatura como $T(t)=\alpha\,\sin\!\big(2\pi\tfrac{t}{24}\big)+\beta$, con mediciones $(t_k,T_k)$, $k=1,\dots,100$. Aunque $T$ es no lineal en $t$, es \emph{lineal} en los parámetros $\alpha,\beta$, así que es un cuadrados mínimos lineal. Matriz de diseño:
\[
A=\begin{pmatrix} \sin(2\pi t_1/24) & 1 \\ \vdots & \vdots \\ \sin(2\pi t_{100}/24) & 1 \end{pmatrix}\in\R^{100\times2},
\quad
x=\begin{pmatrix}\alpha\\\beta\end{pmatrix},
\quad
b=\begin{pmatrix} T_1\\\vdots\\T_{100}\end{pmatrix}.
\]
Se busca $\min_x\norm{Ax-b}_2$ y se resuelve por ecuaciones normales $A^TAx=A^Tb$ (un sistema $2\times2$) o por $x=A^{+}b$. Lo importante de examen: reconocer que un modelo no lineal en $t$ pero lineal en los parámetros se resuelve como MMCC lineal.
\end{ejemplo}

\begin{truco}
\begin{itemize}
\item \textbf{Modelo lineal en los parámetros $\ne$ modelo lineal en $x$.} Ajustes con $x^2$, $x^3$, $\cos(x)$, $\sin(2\pi t/24)$ son todos MMCC \emph{lineales}: una columna por término.
\item Modelos multiplicativos como $V=p_1 x^{p_2}$ se \textbf{linealizan con logaritmo}: $\ln V=\ln p_1+p_2\ln x$, y ahí sí es lineal en $\ln p_1$ y $p_2$.
\item \textbf{Existencia de Cholesky}: $A^TA$ admite Cholesky $\iff$ es definida positiva $\iff$ las columnas de $A$ son LI. Si son dependientes, usá SVD ($x=A^+b$).
\item Cuidado con el condicionamiento: las ecuaciones normales elevan $\kappa_2$ al cuadrado. Si el problema está mal condicionado, preferí QR o SVD.
\item El residuo óptimo $b-Ax^\star$ es ortogonal a $\Imag(A)$: sirve como verificación ($A^T(b-Ax^\star)=0$).
\end{itemize}
\end{truco}
